% Nouveau EGA V (cote 25): archive cote 25, site page 181 (PDF page 182).
% typescript. Draft compared: none.
% AI-transcribed by Claude Opus 5.5 (claude-opus-5-5, adaptive thinking, effort high), Message Batches API, batch ega25b (msgbatch_01A7n5YuXubUPNK1zsepi8RG), 2026-11-13; not yet checked by eye.
\scanpage{25/181}
\mspage{8}

d'une somme directe de formes $Q_1 \oplus Q_2$ , $Q_1$ de rang 1 et $Q_2$ de rang pair.
En vertu de 16.8. on voit que $Q_1,Q_2$ sont de discriminant corrigé
$\neq 0$, en vertu de ce qui précède il s'ensuit que $Q_2$ est isomorphe à la
forme standart, et d'autre part $Q_1$ est non nulle, donc isomorphe à la
forme standart de degré 1, donc Q est isomorphe à la forme standart, ce
qui prouve 16.11.1.

Il reste à prouver l'équivalence de (i) et (ii) dans 16.9. On peut
évidemment supposer que S est le spectre d'un corps alg clos, \inserted[typed]{et on est} \xout{d'autre part}
ramené à prouver dans ce cas \deleted{\unsure{que}} \inserted[typed]{Lemme 16.11.2.} Q lisse \hand{$\Longleftrightarrow$} Q standart.
En vertu de 16.2. on peut supposer que k est de car 2 et E est de rang
impair. \xout{\unsure{On sait déjà que (ii) implique que}} \inserted[typed]{Lorsque} Q est isomorphe à la forme
standart, \xout{\illegible{3}} on vérifie aussitôt, grâce à 16.3. , qu'elle
est lisse (le noyau de $\overset{\hand{\vee}}{E}$ \hand{$\longrightarrow$}Ev est de dimension 1, et Q n'est pas identi-%
quement nulle sur ce noyau, donc n'y admet que le zéro trivial); \xout{donc}
\inserted{D'autre part,} \xout{\unsure{(ii) $\Rightarrow$ (i). Or}} toujours par 16.3. , \inserted{Q lisse} \deleted{(ii)} implique que \xout{le noyau} la restric-%
tion de Q au noyau N de Q n'admet quenle zéro trivial, ce qui évidemment
\inserted{signifie} \xout{implique} que N est de dim 1 et que $Q|N \neq 0$. Or prenant un supplémentaire
de N, on voit aussitôt que cela implique que Q est isomorphe à la forme
standar\over{t}{d}\hand{/} Cela \inserted[typed]{prouve 16.11.2. et} achève la démonstration de 16.9.

NB On aurait pû donner 16.11.2. dès le début en corollaire à
16.2. ou 16.3. ( qui auraient dû s'interchanger); ainsi la partie du n\textsuperscript{o}
indépendante du ``discriminant corrigé'' serait bloquée au début du n\textsuperscript{o}.
D'ailleurs tu dois mieux savoir que moi dans quelle mesure la notion
de discriminant corrigé et 16.9. étaient connus, pour en créditer qui de
droit. Peut-être y a-t\hand{\unsure{-}}il déjà une autre terminologie reçue ? Je remarque
d'ailleurs que c'est le discriminant \uline{\over{n}{c}orrigé} qui mérité vraiment le nom
de discriminant, c'est lui qui se généralise en le discriminant d'une
forme quelconque (cf N\textsuperscript{o} suivant). Le discriminant d'une forme quadratique
(dans la terminologie adoptée ici, et qui, si je ne me trompe, est la

\margin{181}
